\documentclass[10pt,a4paper]{amsart}
\usepackage[margin=19mm]{geometry}
\usepackage{amsmath,amssymb,mathtools}
\usepackage[scr=rsfs,cal=euler]{mathalfa}
\usepackage{xfrac}
\usepackage[unicode,hidelinks,bookmarks=false]{hyperref}
\newcommand{\ie}{i.e.\ }
\newcommand{\cf}{cf.\ }
\newcommand{\resp}{respectively\ }
\newcommand{\Subsection}{\subsection}
\providecommand{\nobreakdash}{\nobreak\hbox{-}}
\setlength{\emergencystretch}{2em}
\multlinegap0pt
\usepackage{amssymb}
\usepackage{mathtools}
\newtheorem{theorem}{Theorem}
\title{Elementary Bounds on Digital Sums of Powers, Factorials, and LCMs}
\author{David G. Radcliffe}
\date{Source: arXiv:2511.15850v2 (CC BY 4.0)}
\begin{document}
\maketitle

\noindent\emph{Source and licence.} Elementary Bounds on Digital Sums of Powers, Factorials, and LCMs. Version arXiv:2511.15850v2 (CC BY 4.0), distributed by arXiv under the Creative Commons Attribution 4.0 International licence, http://creativecommons.org/licenses/by/4.0/, as recorded in the arXiv licence metadata for this exact version and shown on the abstract page https://arxiv.org/abs/2511.15850v2. Full licence notice: \texttt{licenses/arxiv-2511.15850-CC-BY-4.0.txt}.\par\smallskip

\noindent\textit{Editorial scope.} Counted unit: the article's theorem thm:multiples-powers-log-bound (source0.tex lines 280--289). The article's complete original proof of it is delivered with it (source0.tex lines 291--334, one \texttt{proof} environment of 44 lines). No mathematics is written, repaired or re-derived: the statement and the proof are the article's own lines, in the article's own notation, and no retained line is substituted or re-flowed.  Retained uncounted as the source-local context the proof needs: lines 262--268.  Nothing was omitted from any retained span, so the package carries no omission record.  No retained line contains a citation, so the wrapper carries no bibliography and no bibliography entry.  No retained line contains a figure environment, an image-inclusion command or drawing code, so no image is delivered and the package declares no asset dependency.  Editorial formatting only, in addition: an AMS \texttt{amsart} preamble that restates the article's own declarations and macros used by the retained lines, namely the environments \texttt{theorem} on separate counters, together with the standard packages those lines need.  The retained environments print in this document as Theorem 1 in order of appearance, whereas the article prints the same items with the numbering of its own full document.  The complete native source of the article as distributed in the arXiv e-print for this exact version is bundled unchanged as \texttt{sources/189477/source0.tex}.
\begin{theorem} \label{thm:multiples-of-powers}
    Let $2 \le a < b$ be integers with $a \mid b$.
    Let $(e_{k})$ be a sequence of integers such that $e_{1} \ge 1$ and
    $a^{e_{k}} > b^{e_{k-1}}$ for all $k \ge 2$.
    Suppose that $N$ is divisible by $a^{e_{k}}$ but not by $b$.
    Then $c_b(N) \ge k$.
\end{theorem}

\begin{theorem} \label{thm:multiples-powers-log-bound}
    Let $2 \le a < b$ be integers with $a \mid b$.
    Suppose that $N$ is divisible by $a^n$ but not $b$. Then
    there exists $C > 0$, depending only on $a$ and $b$, such that
    \[
        c_b(N) > C \log n
    \]
    for $n$ sufficiently large.
    Moreover, any constant $0 < C < (\log(\log(b)/\log(a)))^{-1}$ is admissible.
\end{theorem}

\begin{proof}
    Let $r = \log(b) / \log(a)$, and
    define $(e_k)$ by $e_1 = 1$ and $e_k = \lceil r e_{k-1} \rceil$
    for $k \ge 2$.
    It is routine to verify that $(e_k)$ satisfies the conditions of
    Theorem~\ref{thm:multiples-of-powers}.

    By the definition of $(e_k)$, we have
    \begin{align*}
        e_2 & < r + 1,             \\
        e_3 & < r^2 + r + 1,
    \end{align*}
    and in general,
    \begin{equation} \label{eq:ek-estimate}
        e_k < \sum_{i=0}^{k-1} r^{i}
        < \frac{r^k}{r - 1}.
    \end{equation}

    Fix an integer $n \ge r/(r-1)$, and let
    \[
        k = \left\lfloor \frac{\log((r-1)n)}{\log r} \right\rfloor.
    \]
    Then \[
        1 \le k \le \frac{\log((r-1) n)}{\log r},
    \]
    which implies that
    \[
        n \ge \frac{r^k}{r-1}.
    \]
    Therefore $n > e_k$ by~\eqref{eq:ek-estimate}, hence
    $c_b(N) \ge k$ by Theorem~\ref{thm:multiples-of-powers}.

    Since
    \[
        \left\lfloor \frac{\log((r-1)n)}{\log r} \right\rfloor
        = \frac{\log n}{\log r} + O(1),
    \]
    choosing any $C < (\log r)^{-1}$ gives
    \[
        \left\lfloor \frac{\log((r-1)n)}{\log r} \right\rfloor
        > C \log n
    \]
    for $n$ sufficiently large, which implies that $c_b(N) > C \log n$.
\end{proof}

\end{document}
